\documentclass{article}
\usepackage[utf8]{inputenc}
\usepackage{amsmath, amsthm, amssymb}
\usepackage{booktabs}
\usepackage{graphicx}
\usepackage{tcolorbox}
\usepackage{hyperref}

\title{ResLocoTransformer: Project Understanding Report}
\author{Gemini CLI Research Agent}
\date{\today}

\begin{document}

\maketitle

\section{Overview}
The goal of this project is to implement and train a transformer-based Reinforcement Learning (RL) policy for the Unitree Go2 quadruped robot using the MuJoCo physics engine. The project emphasizes visual-proprioceptive fusion and explores novel transformer architectures like the \texttt{LocoAttnResTransformer}.

\section{Environment: UnitreeMujocoGymEnv}
The environment is implemented in \texttt{envs/mujoco_env.py}.

\paragraph{Robot and Physics}
\begin{itemize}
    \item \textbf{Robot}: Unitree Go2.
    \item \textbf{Simulator}: MuJoCo 3.x.
    \item \textbf{Control}: Joint position targets (12-dim), converted to torques via PD control ($K_p=40, K_d=0.5$).
\end{itemize}

\paragraph{Observations}
The observation is a concatenated vector of size $84 + 16384 = 16468$.
\begin{itemize}
    \item \textbf{State (84-dim)}: Proprioceptive history (length 3).
    \begin{itemize}
        \item IMU: [roll, pitch, gyro_x, gyro_y] (4 dims $\times$ 3 steps = 12).
        \item Joint Angles: 12 joints $\times$ 3 steps = 36.
        \item Last Actions: 12 joints $\times$ 3 steps = 36.
    \end{itemize}
    \item \textbf{Depth (16384-dim)}: 4 stacked $64 \times 64$ depth frames.
    \begin{itemize}
        \item Rendered from a body-attached forward-facing camera.
        \item Log-transformed: $d' = \sqrt{\log(d+1)}$.
        \item Optionally z-scored.
    \end{itemize}
\end{itemize}

\paragraph{Reward Function}
The primary reward terms include:
\begin{itemize}
    \item \textbf{Velocity Tracking}: Exponential reward for matching target forward velocity.
    \item \textbf{Yaw Rate Tracking}: Exponential reward for matching target angular velocity.
    \item \textbf{Survival Bonus}: Constant positive reward per step.
    \item \textbf{Penalties}: Torque, action rate (smoothness), orientation (roll/pitch), height deviation, and lateral velocity/yaw rate.
\end{itemize}

\section{Model Architecture}
The architectures are defined in \texttt{models/nets.py} and \texttt{torchrl/networks/}.

\paragraph{LocoTransformerEncoder}
\begin{itemize}
    \item \textbf{Visual}: 3-layer CNN (NatureEncoder) $\to$ 16 tokens ($4 \times 4$ grid) of size \texttt{token_dim}.
    \item \textbf{State}: MLP $\to$ 1 token of size \texttt{token_dim}.
    \item Tokens are concatenated into a sequence of length 17.
\end{itemize}

\paragraph{LocoTransformer}
Standard Transformer layers followed by mean/max pooling of visual tokens and concatenation with the state token.

\paragraph{LocoAttnResTransformer}
Uses \textbf{Attention Residuals} (\texttt{AttnResLayer}). Standard additive residuals are replaced with depth-wise attention over all preceding layer outputs.
\[ x_{l+1} = \text{SoftmaxAttention}(q=f(x_l), K=H_l, V=H_l) \]
where $H_l = [x_0, x_1, \dots, x_l]$.

\section{Training Pipeline}
\begin{itemize}
    \item \textbf{Algorithm}: Proximal Policy Optimization (PPO).
    \item \textbf{Framework}: Custom \texttt{torchrl} library.
    \item \textbf{Normalization}: \texttt{NormObsWithImg} normalizes the state portion only.
    \item \textbf{Logging}: Weights and Biases (WandB) and local TensorBoard logs.
\end{itemize}

\section{Experiment Summary Table}
\begin{table}[h!]
\centering
\setlength{\tabcolsep}{8pt}
\renewcommand{\arraystretch}{1.2}
\begin{tabular}{lllrrc}
\toprule
ID & Date & Description & Commit & Metric & Status \\
\midrule
E000 & 2026-06-08 & Baseline Understanding & N/A & N/A & \textbf{Complete} \\
E001 & 2026-06-08 & Baseline Play Test & b8fc464 & N/A & \textbf{Failed (Arch Incompatible)} \\
\bottomrule
\end{tabular}
\caption{Summary of experiments.}
\end{table}

\section{Analysis of Baseline}
Attempting to run \texttt{scripts/play.py} on the baseline checkpoint (\texttt{00\_baseline}) resulted in a architecture mismatch. The current code assumes \texttt{LocoAttnResTransformer}, while the checkpoint seems to use a different (likely standard) transformer architecture.
\end{document}
